\chapter{기계는 어떻게 배우는가}
% full version: chapters/06_dofa.tex

지금까지의 모든 회로에서는 연결의 세기를 엔지니어가 골랐다. 살아 있는 뇌에는 엔지니어가 없다. 연결은 일하는 도중에 스스로 조정되어야 하고, 그 결과 기계가 쓸모 있는 일을 해야 한다. 이것이 바로 학습이다.

핵심 제약은 매우 엄격하다. 각 연결은 스스로를 바꾸며, 알 수 있는 것은 가까이에서 일어나는 일뿐이다. 송신자가 일하는지, 수신자가 일하는지, 어떤 물질이 자기에게 닿는지. 누구도 연결에 개인별 수정값을 보내 줄 수 없다. 이 때문에 인공 신경망 학습의 주된 방법은 쓸 수 없다. 거기서는 오차가 네트워크를 거꾸로 거슬러 올라가, 연결마다 제 몫의 수정값을 받는다.

\section*{함께 일하면 이어진다}

가장 단순한 규칙은 이렇다. 송신자와 수신자가 함께 일하면 그 사이의 연결이 강해진다. 이 규칙은 국소적이고 시계가 필요 없다. 여기에 약간의 망각을 더하면, 이 규칙은 신호의 흐름에서 가장 중요한 것, 즉 가장 크게 변하는 것을 찾아낸다. 딸깍의 시각을 보는 변형도 있다. 송신자가 수신자\emph{보다 먼저} 딸깍했다면, 즉 원인이었을 수 있다면 연결이 강해지고, 나중에 딸깍했다면 약해진다.

그러나 이런 규칙에는 한계가 하나 있다. 이 규칙은 \emph{자주} 일어나는 것을 가르치지, \emph{쓸모 있는} 것을 가르치지 않는다. 이웃만 보는 연결은 어떤 행동이 주스를 가져왔는지 고통을 가져왔는지 알지 못한다.

\section*{세 번째 신호}

그것을 알려 주는 것이 도파민, 곧 “예상보다 좋다” 방송국이다. 연결은 세 가지가 겹칠 때만 바뀌게 하자. 송신자가 일했고, 수신자가 일했고, 도파민이 왔을 때이다.

\pencilfig{6}{\pencilart{6}}{두 뉴런이 함께 일했고 세 번째 신호 “기대보다 잘됐다”가 오면 연결이 강해진다.}

그런데 어려움이 있다. 보상은 곧바로 오지 않고 1, 2초 뒤에 온다. 그때쯤이면 송신자도 수신자도 벌써 조용해졌다. 연결에는 자기가 참여했다는 기억이 필요하다. 그리고 그 기억은 있다. 함께 일한 흔적이 연결에 표시로 남는다. 1, 2초 뒤면 저절로 떨어지는 스티커와 같다. 스티커가 붙어 있는 동안 도파민이 오면 연결이 바뀐다. 오지 않으면 표시는 흔적 없이 사라진다. 이것으로 주소 지정 문제가 풀린다. 도파민은 모두에게 하나지만, 바뀌는 것은 스티커가 붙은 연결뿐이다.

\pencilfig{106}{%
\foreach \y/\d/\t in {0/0.9/{제때: 연결 강화},-1.3/3.3/{늦음: 변화 없음}}{
  \draw[pen] (0,\y) -- (4.6,\y);
  \foreach \s in {0.25,0.33}{\draw[pcBlue, line width=0.8pt] (\s,\y) -- (\s,\y+0.3);}
  \fill[pcAmber, opacity=0.25] (0.35,\y) -- plot[domain=0.35:4.5, samples=30] (\x,{\y+0.8*exp(-(\x-0.35)/0.8)}) -- (4.5,\y) -- cycle;
  \draw[penlite=pcAmber] plot[domain=0.35:4.5, samples=30] (\x,{\y+0.8*exp(-(\x-0.35)/0.8)});
  \draw[pcAmber!80!black, line width=0.9pt] (\d-0.1,\y+0.95) -- (\d+0.07,\y+0.62) -- (\d-0.07,\y+0.6) -- (\d+0.1,\y+0.22);
  \draw[arr=pcAmber!80!black] (\d+0.09,\y+0.27) -- (\d+0.11,\y+0.12);
  \node[lbl, anchor=west] at (4.7,\y+0.3) {\t};}
\node[lbl, text=pcBlue, anchor=south west] at (0.1,0.85) {함께 일했다};
\node[lbl, text=pcAmber!80!black, anchor=south] at (2.3,0.35) {표시가 사라진다};
\foreach \x/\l in {0.35/0,1.95/1,3.55/2}{\draw[pcGray, line width=0.5pt] (\x,-1.3) -- (\x,-1.38); \node[lbl, anchor=north] at (\x,-1.38) {\l~s};}
}{함께 일한 흔적은 연결에 표시로 남고, 이 표시는 1, 2초 만에 사라진다. 도파민은 표시가 있는 연결만 바꾼다.}

\section*{탐색으로서의 잡음}

이 규칙은 왜 아무 데로나가 아니라 더 나은 쪽으로 이끄는가? 눈을 가린 채 언덕 꼭대기를 찾는다고 해 보자. 아무렇게나 한 걸음 내딛으면 누군가 “올라갔다!” 또는 “내려갔다!” 하고 외친다. “올라갔다”면 그 걸음을 굳힌다. 지도를 한 번도 보지 않고도 한 걸음씩 올라간다. 뇌도 똑같이 한다. 뉴런들은 무작위로 조금씩 흔들리고, 도파민은 나아졌는지 알려 주며, 참여한 연결은 유리한 쪽으로 옮겨 간다. 메시지 전달을 방해하던 잡음이 여기서는 탐색이 된다.

이제 도파민이 왜 보상이 아니라 \emph{뜻밖의 일}을 알리는지 이해된다. 도파민이 성공할 때마다 그냥 “주스!”라고 외친다면, 네트워크가 하던 모든 일을 옳든 그르든 함께 굳혀 버릴 것이다. 우리 모델에서 이런 네트워크는 실제로 연결을 수천 배까지 키워 버리고, 때로는 더 나쁜 선택에 영원히 갇힌다. 기대를 빼야 학습이 공정해진다.

\section*{전선 하나의 값}

브로드캐스트 학습에는 대가가 있다. 신호는 모두에게 하나이고, 그 안에는 결과에 영향을 주는 모든 뉴런의 잡음이 섞여 있다. 네트워크가 클수록 각 연결이 제 몫을 가려내는 데 오래 걸린다. 그래서 도파민은 행동을 고르는 비교적 작은 회로만 가르치고, 거대한 피질은 주로 국소 규칙으로 배우는 것으로 보인다.

\section*{뜻밖의 정도는 누가 계산하나}

질문이 남는다. 도파민 뉴런은 “예상보다 좋다”를 어떻게 계산하는가? 비평가가 필요하다. 앞으로 좋은 일이 얼마나 더 남았는지 늘 평가하는 뉴런 무리이다. 뜻밖의 정도는 받은 보상에 기대가 얼마나 뛰어올랐는지를 더한 것이다. 전구가 켜지면 기대가 뛰어오르고, 분출이 일어난다. 약속된 주스가 오면 보상은 있지만 기대가 바로 그 순간 소진되므로 뜻밖의 일은 없다. 주스가 오지 않으면 기대가 헛되이 무너지고, 침묵이 생긴다. 첫 장의 세 경우가 모두 저절로 나온다. 도파민이 바로 이렇게 행동한다는 것은 측정된 사실이다. 뇌가 이 뜻밖의 정도를 정확히 어떻게 계산하는지는 아직 가설이다.

\fullversion{6}
